\documentclass[11pt]{beamer}
\usetheme{cambridgeUS}
\usepackage[utf8]{inputenc}
\usepackage[ngerman]{babel}
\usepackage{amsmath}
\usepackage{amsfonts}
\usepackage{amssymb}
\usepackage{graphicx}
\usepackage{listings}
\usepackage{tikz}
\usetikzlibrary{automata,positioning}

\author{Maria Kosche}
\title{Latex Einführung}
\subtitle{Basierend auf einem Vortrag von Janina Schmidt}
%\setbeamercovered{transparent} 
%\setbeamertemplate{navigation symbols}{} 
%\logo{} 
%\institute{} 
%\date{} 
%\subject{} 
\begin{document}

\begin{frame}
\titlepage
\end{frame}

\section{Warum Latex?}
\begin{frame}{Warum Latex?}
\begin{itemize}
	\item Professionelles Layout
	\item Einfache Umsetzung von mathematischem Zeichensatz
	\item Verwendung für Bachelor- und Masterthesis
	\item Unterstützung auf allen Plattformen
\end{itemize}
\end{frame}

\section{Welcher Editor?}
\begin{frame}{Welcher Editor?}
\begin{itemize}
	\item<1-> Sharelatex als online Lösung der GWDG
	\url{https://sharelatex.gwdg.de}
	\item<2-> Kile auf den Rechnern des CIP-Pools
	
	\texttt{sudo apt-get install kile}
	\item<4-> texmaker 
	
	\texttt{sudo apt-get install texmaker}
	\item<5-> TeXstudio
	\url{https://www.texstudio.org/}
	\item<6-> Hilfe und Dokumentation unter \url{https://en.wikibooks.org/wiki/LaTeX}
\end{itemize}
\end{frame}

\section{Minimalbeispiel}
\begin{frame}[fragile]
\frametitle{Minimalbeispiel}

\begin{itemize}
	\item Grundsätzlich unterteilen wir ein \LaTeX-Dokument in Preamble und Body
	\item Preamble stellt das Fundament dar (Dokumentenklasse, Spracheinstellungen, zu ladende Pakete, $\ldots$)
	\item Body enthält den zu setzenden Text:
	$\rightarrow$ \verb|\begin{document} ... \end{document}|
\end{itemize}

\begin{verbatim}
\documentclass[11pt]{article}

% Kommentare in LaTeX beginnen mit einem %-Zeichen
% In der Preamble werden Dokumenteinstellungen festgelegt und Pakete geladen

\begin{document}
% Im Dokumentenkörper wird der eigentliche Text gesetzt.
Hello, World.
\end{document}
\end{verbatim}
\end{frame}


\section{Dokumentklasssen}
\begin{frame}{Dokumentklassen}
\begin{itemize}
\item article
\item beamer
\item letter
\item book
\item tikzposter
\item Definition eigener Dokumentklassen möglich
\end{itemize}
\end{frame}

\section{Dokumentstruktur}
\begin{frame}[fragile]
\frametitle{\secname}
\begin{itemize}
	\item<1-> Aufteilen des Dokuments in logische Abschnitte
	\item<2-> \begin{verbatim}
	\section{name} \subsection \subsubsection{} 
	\paragraph{•}
	\end{verbatim}
	\item<3-> Standardmäßig nummeriert, wenn nicht gewünscht \verb|\section*{}|
	\item<4-> Automatisch erstelltes und sich selbst anpassendes Inhaltsverzeichnis mit \textbackslash tableofcontents
	\item<5-> Im erstellten pdf-Dokument sind Seiten über das Inhaltsverzeichnis verlinkt.
\end{itemize}
\end{frame}

\begin{frame}[fragile]
	\frametitle{\secname\ II}
	\begin{itemize}
		\item<1-> Sektionen, Kapitel usw. für die logische Strukturierung des Textes
		\item<2-> Sogenannte Umgebungen für die Strukturierung des Textsatzes
		\item<3-> Umgebungen legen fest, wie Text gesetzt werden
		\item<4-> Beispiel einer Tabellenumgebung: \verb|\begin{tabular} ... \end{tabular}|
		\item<5-> Weitere Beispiele und Verwendungshinweise in diesem Vortrag
	\end{itemize}
\end{frame}


\section{Textformatierung -- Schrift}
\begin{frame}[fragile]
\frametitle{\secname}
\begin{itemize}
	\item<1-> Reiner Text, der mit Latex gesetzt wird, kann mittels verschiedener Befehle formatiert werden
	\item<2-> Für \textbf{Fettschrift} kann beispielsweise \verb|\textbf{<text>}| verwendet werden
	\item<3-> Schriftgröße wird grundsätzlich mit der Dokumentenklasse oder in der Preamble ausgewählt (z.B. \verb|\documentclass[11pt]{article}|)
	\item<4-> Auch die {\huge temporäre} {\tiny Veränderung} der Schriftgröße ist möglich mit Befehlen wie \verb|\large| oder \verb|\tiny|
	\item<5-> \url{https://en.wikibooks.org/wiki/LaTeX/Fonts#Font_styles}
\end{itemize}
\end{frame}

\section{Gruppen}
\begin{frame}[fragile]
	\frametitle{\secname}
	\begin{itemize}
		\item<1-> Gruppen: Text mit bestimmten Eigenschaften gruppieren
		\item<2-> Formatierung und Definitionen verlieren bei Gruppenende ihre Wirkung
		\item<3-> \begin{verbatim}
			\begingroup
			\itshape
			Dies ist kursiver Text in einer Gruppe.
			\endgroup
			Nach der Gruppe wird die Formatierung zurückgesetzt.
		\end{verbatim}
		$$\Downarrow$$
		\begingroup
		\itshape
		Dies ist kursiver Text in einer Gruppe.
		\endgroup
		Nach der Gruppe wird die Formatierung zurückgesetzt.
	\end{itemize}
\end{frame}

\section{Aufzählungen}
\begin{frame}[fragile]
\frametitle{Aufzählungen}
\begin{itemize}
\item Um eine Aufzählung zu bekommen kann die \textbackslash itemize Umgebung verwendet werden
\item Eine Aufzählung sieht dann so aus:
\begin{verbatim}
\begin{itemize}
	\item article
	\item book
	\item tikzposter
\end{itemize}
\end{verbatim}
\item[-] Mit \textbackslash \verb|item[-]| können statt Bulletpoints Spiegelstriche verwendet werden
\item[$\gamma$] mit \verb|item[$CHARACTER$]| auch jedes beliebige andere Zeichen
\end{itemize}

\end{frame}

\begin{frame}{Aufzählungen II}
\begin{enumerate}[1]
\item Mit \textbackslash begin\{enumerate\} können geordnete Listen erstellt werden
\item \begin{enumerate}[a]
		\item Zusätzliche Listenebenen sind auch möglich
		\item \ldots
	\end{enumerate}
\item \ldots
\end{enumerate}
\end{frame}

\section{Matheformeln}
\begin{frame}[fragile]
\frametitle{Matheformeln I}
\begin{itemize}
\item<1-> Möchte man Mathematische Ausdrücke im Text verwenden, so muss man den gewünschten Ausdruck in \$ \$ schreiben

Beispiel:
\verb|$i \in \mathbb{N}$| wird zu $i \in \mathbb{N}$
\item<2-> Komplizierte oder größere Formeln sollten in einer Formelumgebung stehen

\$\$ ... \$\$ oder \textbackslash begin \{equation\}

Beispiel:
\begin{verbatim}
	$$ \sum_{i=2}^{5t^{2}} \frac{1}{2 \pi i} $$
\end{verbatim}
wird zu
$$ \sum_{i=2}^{5t^{2}} \frac{1}{2 \pi i} $$
\end{itemize}
\end{frame}

\begin{frame}[fragile]
\frametitle{Matheformeln II}
\begin{itemize}
\item<1-> Griechische Buchstaben werden mit \textbackslash NAME geschrieben, z.B. \textbackslash gamma oder \textbackslash Gamma gibt $\gamma$ und $\Gamma$ 
\item<2-> Brüche werden mit \textbackslash frac\{ZÄHLER\}\{NENNER\} dargestellt
\item<3-> Beliebige Dinge bekommen einen Index durch \_\{ \} und einen Exponenten durch $\hat{}$ \{ \}
\item<4-> Beispiel $\alpha_1^{2x}$ durch \begin{verbatim}
$\alpha_1^{2x}$
\end{verbatim}
\item<5-> Allquantor $\forall$ \textbackslash forall, Existenzquantor $\exists $ \textbackslash exists, $\in$ \textbackslash in, $\notin$ \textbackslash notin
\end{itemize}
\end{frame}

\begin{frame}[fragile]
\frametitle{Matheformeln III}
\begin{itemize}
\item<1-> Summen $\sum$ können durch \textbackslash sum erzeugt werden
\item<2-> Integral $\int$ durch \textbackslash int
\item<3-> Klammern, die durch \textbackslash left( und \textbackslash right) erstellt werden, passen sich an die Größe der Formel an
\item<4-> $\mathbb{N}$, $\mathbb{Q}$ durch \textbackslash mathbb \{BUCHSTABE\}
\item<5-> \url{https://en.wikibooks.org/wiki/LaTeX/Mathematics}
\end{itemize}
\end{frame}

\begin{frame}[fragile]
\frametitle{Matheformeln IV}
\begin{itemize}
\item<1-> Darstellung Matrix mit Array Umgebung $$A = \left( \begin{array}{cc}
a_{11} & a_{12} \\
a_{21} & a_{22} \\
\end{array} \right)$$
\item<2-> \begin{verbatim}
$$ \left( \begin{array}{cc}
a_{11} & a_{12} \\
a_{21} & a_{22} \\
\end{array} \right)$$
\end{verbatim}
\end{itemize}
\end{frame}

\begin{frame}[fragile]
\frametitle{Matheformeln V}
\begin{itemize}
\item Für Termumformungen möchte man untereinder angeordnete Formeln, dies ist mit \verb|\begin{align} ... \end{align}| möglich
\begin{align*}
x^{2} + y^{2} &= 9 \\
x^{2} &= 9 - y^{2}\\
\end{align*}

\begin{verbatim}
\begin{align*}
x^{2} + y^{2} &= 9 \\
x^{2} &= 9 - y^{2}\\
\end{align*}
\end{verbatim}
\end{itemize}
\end{frame}

\section{Tabellen}
\begin{frame}[fragile]
\frametitle{Tabellen I}
\begin{itemize}
\item Einfache Tabellen  können mit einer tabular Umgebung in einer table Umgebung erstellt werden

\begin{table}[]
	\begin{tabular}{l|lll}
		\multicolumn{4}{c}{\textbf{Personenübersicht}} \vspace{0.5cm} \\
		Name&  Alter&  Studienfach&  Hobby \\
		\hline
		Jean&  117&  Irgendwas mit Medien&  Netflix \\
		Jaque&  23&  Informatik& \LaTeX
	\end{tabular}
\end{table}

\footnotesize \begin{verbatim}
\begin{table}[]
	\begin{tabular}{l|lll}
		\multicolumn{4}{c}{\textbf{Personenübersicht}} \vspace{0.5cm} \\
		Name&  Alter&  Studienfach&  Hobby \\
		\hline
		Jean&  117&  Irgendwas mit Medien&  Netflix \\
		Jaque&  23&  Informatik& \LaTeX
	\end{tabular}
\end{table}
\end{verbatim}
\end{itemize}
\end{frame}

\begin{frame}[fragile]
\frametitle{Tabellen II}
\begin{itemize}
\item<1-> $l$, $c$ und $r$ in dem \textbackslash begin\{tabular\}\{llll\} Befehl geben für jede Spalte an, wie sie ausgerichtet sein soll
\item<2-> $|$ und \textbackslash hline für gerahmte Spalten/Zeilen
\item<3-> Beschriftung und Mehrfachzellen sind möglich mit den Befehlen \verb|\caption{}|
und \verb|\multicolumn{col}{pos}{text}|
\end{itemize}
\pause\pause\pause
\begin{table}[]
\caption{Tabelle}
\begin{tabular}{|c|c|c|c|}
 Name&  Alter&  Studienfach&  Hobby \\
 \hline
 \multicolumn{2}{|c|}{ 2}&  2&  2
\end{tabular}
\end{table}

\footnotesize \begin{verbatim}
\begin{table}[]
\caption{Tabelle}
\begin{tabular}{|c|c|c|c|}
 Name&  Alter&  Studienfach&  Hobby \\
 \hline
 \multicolumn{2}{|c|}{ 2}&  2&  2
\end{tabular}
\end{table}
\end{verbatim}
\end{frame}

\section{Figures}
\begin{frame}[fragile]
\frametitle{\secname\ I}
\begin{itemize}
	\item Verwenden von Graphiken mit dem graphicx Paket \verb|\usepackage{graphicx}|
\end{itemize}

\begin{figure}[h]
	\centering
	\includegraphics[width=0.3\linewidth]{test.pdf}
	\caption{Dies ist nur ein Testbild}
\end{figure}

\begin{verbatim}
\begin{figure}[h]
	\centering
	\includegraphics[width=0.3\linewidth]{test.pdf}
	\caption{Dies ist nur ein Testbild}
\end{figure}
\end{verbatim}
\end{frame}

\begin{frame}[fragile]
\frametitle{\secname\ II}
	\begin{itemize}
		\item Verschiedene float specifier, um Graphiken auf der Seite zu platzieren.
		
		[h] (for here), [t] (top of the page), [b] (bottom of the page)
		\item Anpassung der Breite oder Höhe einer Graphik mit den \verb|width|- und \verb|height|-Optionen
		\item Bildbeschriftung mit dem \verb|caption{}|-Befehl
		\item \url{https://en.wikibooks.org/wiki/LaTeX/Floats,_Figures_and_Captions}
	\end{itemize}
\end{frame}

\section{Code}
\begin{frame}[fragile]
\frametitle{Code}
\begin{itemize}
\item<1-> \textbackslash usepackage\{listings\}
\item<2-> \footnotesize \lstset{language=Python} 
\begin{lstlisting} 
def dct1(a):
	y = np.empty(2*n, dtype=np.float64)
	y[0] = math.sqrt(2)*a[0]
	y[n+1:2*n] = a[n-1:0:-1]
	return erg
	\end{lstlisting}
\item<3-> \begin{verbatim}
\lstset{language=Python}
\begin{lstlisting} 
def dct1(a):
	y = np.empty(2*n, dtype=np.float64)
	y[0] = math.sqrt(2)*a[0]
	y[n+1:2*n] = a[n-1:0:-1]
	return erg
	\end{lstlisting}
\end{verbatim}
\end{itemize}
\end{frame}

\section{Graphen}
\begin{frame}[fragile]
\frametitle{Graphen I}
\begin{itemize}
\item Mit dem tikz Paket kann man Graphiken zeichnen \textbackslash usepackage\{tikz\}
\begin{tikzpicture}[shorten >=1pt,node distance=2cm,on grid,auto] 
   \node[state,initial] (q_0)   {$q_0$}; 
   \node[state] (q_1) [above right=of q_0] {$q_1$}; 
   \node[state] (q_2) [below right=of q_0] {$q_2$}; 
   \node[state,accepting](q_3) [below right=of q_1] {$q_3$};
    \path[->] 
    (q_0) edge  node {0} (q_1)
          edge  node [swap] {1} (q_2)
    (q_1) edge  node  {1} (q_3)
          edge [loop above] node {0} ()
    (q_2) edge  node [swap] {0} (q_3) 
          edge [loop below] node {1} ();
\end{tikzpicture}
\end{itemize}
\end{frame}

\begin{frame}[fragile]
\frametitle{Graphen II}
 \scriptsize \begin{verbatim}
\begin{tikzpicture}[shorten >=1pt,node distance=2cm,on grid,auto] 
   \node[state,initial] (q_0)   {$q_0$}; 
   \node[state] (q_1) [above right=of q_0] {$q_1$}; 
   \node[state] (q_2) [below right=of q_0] {$q_2$}; 
   \node[state,accepting](q_3) [below right=of q_1] {$q_3$};
    \path[->] 
    (q_0) edge  node {0} (q_1)
          edge  node [swap] {1} (q_2)
    (q_1) edge  node  {1} (q_3)
          edge [loop above] node {0} ()
    (q_2) edge  node [swap] {0} (q_3) 
          edge [loop below] node {1} ();
\end{tikzpicture}
\end{verbatim}
\begin{itemize}
\item \url{https://en.wikibooks.org/wiki/LaTeX/PGF/TikZ}
\item Verwendung von Zeichentools ebenfalls möglich (z.B. Ipe \url{http://ipe.otfried.org/})
\end{itemize}
\end{frame}

%\section{Aufgaben}
%\begin{frame}{Aufgaben}
%\begin{enumerate}
%\item Erstelle ein Dokument dass aus mindestens zwei Kapiteln besteht
%\item Erstelle ein Inhaltsverzeichnis für das Dokument
%\item Erstelle in dem Dokument einen Text, der an einem Beispiel die Verwendung der PQ-Formel erklärt
%\item Erstelle in dem Dokument eine Tabelle, die den Namen, das Alter und ein Hobby der anderen Workshopteilnehmer*innen enthält
%\end{enumerate}
%\end{frame}
\end{document}